\section{Unloosening the Knots}

\begin{quotationx}
The Buddha took a silken handkerchief, tied a knot in it, held it up and asked Ananda: ``What is this?'' Ananda replied: ``A silk handkerchief with a knot.'' The Buddha tied a second knot, then a third, and so on up to six knots. After each knot he asked Ananda the same question and got the same response.

The Buddha then said: ``When I tied the first knot, you called it a knot, and so on.'' Ananda became puzzled and exclaimed: ``Whether you tie a single knot or a hundred, they remain knots, though the handkerchief is made of variously coloured pieces of silk.''

The Buddha admitted this and pointed out that although the piece of silk was one, and all the knots were knots, the difference lay in the order they were tied.

The Buddha asked how these knots could be untied. He pulled at the knots, but the knots became tighter. Ananda replied: ``I would first find out how the knots were tied. For he who knows the origin of things knows also their dissolution. But can all the knots be untied at the same time?''

The Buddha replied: ``Not at all! Since the knots were tied one after another in a certain order, we cannot untie them, unless we follow the reverse order.''
\end{quotationx}


\textbf{Lama Angarika Govinda} relates these knots to the chakras, which are prior to any activities of the sense organs. He claims that this is why the path must start at the lowest chakras rather than the higher. In other words, we have to reverse the descent of the spirit into matter (``the coagulation of consciousness into a state of materiality'') by untying the knots one by one in the reverse order in which we created them.

In Man and his Becoming, \textbf{Rene Guenon} describes this descent from Purusha (spirit) to Prakriti (prime matter). Ontologically and logically, the higher states are prior, but the lower states are created first in a temporal sequence. It is useful to try to understand this in the terms of the Western tradition, with its notion of levels of the soul\footnote{\url{https://www.gornahoor.net/?p=3390}}. In this understanding, the higher levels of the soul are to dominate the lower. Knots in the lower soul prevent the free flow of energy from the higher to the lower.

\paragraph{Attraction and Aversion}

The nutritive soul reacts through attraction and aversion. For example, a plant will bend toward the sun, although desert plants may turn away. In human consciousness, this is experienced as likes and dislikes. These take hold in our minds prior to conscious awareness and are thus automatic processes. As in the case for plants, these attractions and aversions may be useful. However, such as we are, they more often create knots in our nutritive souls, since they subject us to, for the most part, arbitrary attractions and aversions, from long forgotten sources, which no longer suit us in a healthy way.

Nevertheless, this low level of awareness is very common. Men make decisions, even important decisions, on the basis of what they like or dislike, similar to tropism in a plant\footnote{\url{https://www.gornahoor.net/?p=8349}}. To unravel this first knot requires the awareness of our likes and dislikes and how they affect us.

\paragraph{Desire}


In the sensitive soul, the analogous process is experienced as desire. When desire is ordered to a man's happiness, it is well functioning. However, knots may appear as weakness (disorder of \emph{thumos\footnote{\url{https://www.gornahoor.net/?p=1956}}) and concupiscence (disorder of
\emph{eros})}. The way these knots manifest can be catalogued and be brought into conscious awareness.

\paragraph{Will}


Since the functions of the intellectual soul, or lower intellect, include knowledge and will, knots appear as ignorance and impotence (or malice). Concerted efforts must be made to gain knowledge and power.

\paragraph{Buddhi and Logos}


The highest, or transpersonal, states do not have an opposite. However, since they transcend the individual states, they are beyond our awareness insofar as we are knotted in an individual state. The first degree of manifestation is Buddhi, or higher intellect. Although it is formless, it still belongs to manifestation. Nonetheless,

\begin{quotex}
Buddhi transcends the domain not only of human individuality, but of every individual state whatsoever... If the Self, or Atma, is the Sun which shines at the centre of the entire being, Buddhi will be the ray directly emanating from this Sun and illuminating in its entirety the particular individual state ... while at the same time linking it to the other individual states of the same being.
\end{quotex}


But what is interesting to us most especially, is the symbolism of the Sun's reflection in the water. As Guenon explains it:

\begin{itemize}


\item The living soul is the image of the sun in water as the reflection of the ``Universal Spirit'' (Atma)

\item The luminous ray confers existence upon this image

\item The water reflecting~the solar light is the symbol of universal passivity, common to all traditions

\item Buddhi consequently derives from Prakriti, of which it is the first production

\item The water represents the potential sum of formal possibilities, that is, the individual

\item The sun's reflection is affected by the undulations and perturbations of the water
\end{itemize}


Note how this corresponds precisely with the Hermetic understanding\footnote{\url{https://www.meditationsonthetarot.com/the-word-is-made-flesh}}. The Buddhi is born from the reflection of the Holy Spirit on the waters. The more the waters are clear and calm, the more the reflection will be pure. The waters are disturbed by eddies and whirlpools, which are the knots of lower consciousness. This birth, of course, is not in a temporal sequence.

All doubts can be removed when Guenon points out that what the Vedantists mean by Buddhi is exactly what the Alexandrian theologians meant by Logos. So, in Western terms, the Hindu doctrine is that the Logos is born from the Holy Spirit and the Virgin.


\flrightit{Posted on 2012-02-19 by Cologero}

\begin{center}* * *\end{center}

\begin{footnotesize}
\begin{sffamily}
\texttt{mleow on 2012-02-20 at 03:39 said:}

Where is that initial qoute taken from featuring Buddha and Ananda? Is it some mahayana source behind it? I can`t recall reading it in the pali canon. Surangama Sutra 

\hfill

\end{sffamily}
\end{footnotesize}

